\section{Agent Roles}
\label{sec:agent_roles}

\noindent The Actor Agent's job is to mutate the configuration state of the Gaussian Splatting reconstruction. The actor holds the current config file containing parameters such as number of Gaussians, densification interval, learning rates for position/opacity/scale/rotation, spherical harmonics degree, pruning thresholds, and camera pose refinement flags and applies edits to it based on structured feedback it receives from the Critic. Instead of doing generation from scratch for each iteration, the Actor applies edits such as "increase the densification interval from 100 to 150 steps. The Actor receives the Critic's structured feedback, reasons about which config parameters are causally responsible for the reported artifact, proposes a change to the config, applies it, triggers a re-train or re-optimize pass via gsplat, and then passes the newly rendered output back to the Critic for the next round. This gradient-free optimization loop over the config space, driven by visual and geometric feedback, is better than a loss function, because many failure modes in 3DGS (e.g. floaters, needle Gaussians, blurry textures, background collapse) are not easily encoded as a differentiable objective function.\\

\newpage
\noindent\title{Actor Agent Prompt:}
\begin{figure}[h]
\centering
\begin{verbatim}
prompt = (
    "You are an expert ML 
    Engineer specializing 
    in 3D Gaussian Splatting.\n"
    "A visual critic has 
    reviewed the latest 
    render and provided this 
    feedback:\n\n"
    f"{state['feedback']}\n\n"
    "Current training config:\n"
    f"{current_config_str}\n"
    "Rewrite the config YAML 
    with adjusted hyperparameters 
    to address the critique.\n"
    "Reply with ONLY valid YAML in 
    the exact same structure. 
    No explanation.\n"
    "training:\n"
    "max_steps: <int>\n"
    "densify_grad_threshold: <float>\n"
    "opacity_cull_threshold: <float>"
    )
\end{verbatim}
\caption{Example Actor Agent Prompt}
\label{fig:actor-agent-prompt}
\end{figure}

\noindent The Critic Agent's job is to look at the rendered 3DGS output and produce structured, actionable feedback that the Actor can use to change specific parameters. The Critic will inspect the scene from multiple angles. It examines close-up views of problematic regions, depth maps to detect geometric inconsistencies, opacity visualizations to catch floaters, and normal maps to identify surface reconstruction failures\\

\noindent For the Critic agent, we propose using a powerful VLM such as GPT-4o or Claude Opus. We take inspiration from frameworks like Feature4X, where gpt-4o act as high-level agents for executing and refining complex scene manipulations. In our architecture, the VLM Critic functions as an expert diagnostic layer. By analyzing a multi-view grid of renders alongside depth and opacity buffers, the agent can pinpoint specific optimization failures. For instance, the Critic might identify that “inward-pointing needle Gaussians on a ceiling” suggest a lack of upward-facing or overhead views. Alternatively, it could recognize that a “foggy” floor surface indicates an insufficient spherical harmonics degree, preventing the model from resolving sharp, view-dependent reflections.\\

\newpage
\noindent\title{Critic Agent Prompt:}
\begin{figure}[h]
\centering
\begin{verbatim}
prompt = (
    "You are an expert Machine 
    Learning Critic specializing 
    in 3D Gaussian Splatting.\n"
    "Analyze the rendered novel viewpoints 
    above alongside these 
    training metrics:\n\n"
    f"{metrics_str}\n\n"
    "Look specifically for:\n"
    "1. Floaters / clouding (foggy 
    artifacts in empty space)\n"
    "2. Blurring (loss of 
    fine surface detail)\n"
    "3. Anisotropic stretching 
    (needle-shaped Gaussians)\n\n"
    "Reply in exactly this format:\n"
    "STATUS: ACCEPTED or REJECTED\n"
    "CRITIQUE: <one sentence assessment>\n"
    "RECOMMENDATION: <specific hyperparameter 
    adjustments if REJECTED, else None>"
    )
\end{verbatim}
\caption{Example Critic Agent Prompt}
\label{fig:critic-agent-prompt}
\end{figure}


\noindent The design constraint is that the Critic's output must be structured enough for the Actor to parse and act on a typed schema rather than a free-text paragraph the Actor has to interpret. This structured output lets us replay the full sequence of Critic observations and Actor edits to understand exactly how the scene evolved across iterations.\\

\noindent\title{Critic Agent Feedback:}
\begin{figure}[h]
\centering
\begin{verbatim}
[Critic]: STATUS: REJECTED
CRITIQUE: The model exhibits catastrophic 
blurring and extreme loss of fine 
detail across the novel views, 
as evidenced by the severely 
low PSNR and SSIM metrics.
RECOMMENDATION: Increase the maximum 
Gaussian count and refine the 
densification parameters (e.g., decrease the 
densification threshold) to counteract 
the severe anisotropic stretching and 
improve spatial resolution.
\end{verbatim}
\caption{Sample Critic Agent Feedback}
\label{fig:critic-agent-prompt}
\end{figure}